% реферат по теме 'Аутентификация и доступ'
\documentclass{SibFU-docs}

% доп. пакеты
\usepackage{graphicx}
\usepackage{float}
\usepackage{hyperref}
\usepackage{caption}
\usepackage{subcaption}
\usepackage{amsmath}
% предотвартим ошибку с \Bbbk перед загрузкой пакета amssymb
\let\Bbbk\relax
\usepackage{amssymb}
\usepackage{booktabs}
\usepackage{tabularx}
\usepackage{longtable}
\usepackage{enumitem}
\usepackage{multirow}

% гиперссылки
\hypersetup{
    colorlinks=true,
    linkcolor=black,
    citecolor=blue,
    urlcolor=blue
}

% библиография
\addbibresource{report.bib}

% титульный лист
\begin{document}
\setlist[itemize]{left=1.3cm, itemsep=1pt, topsep=2pt} % для выравнивания списков
\setlist[enumerate]{left=1.3cm, itemsep=1pt, topsep=2pt}
\renewcommand{\contentsname}{}
\mycovertitle
{Гуманитарный институт}
{Кафедра прикладной информатики в искусстве и интерактивном медиа}
{Информационные технологии и информационные системы}
{ст. преподаватель И.~Р.~Нигматуллин}
{Логинова Ю.В., Захаров И.М.}
 
\begin{center}
\bfseries РЕФЕРАТ
\end{center}
\noindent

Реферат по теме: «Информационные технологии и информационные системы» содержит 23 страниц и 15 использованных источников.

Ключевые слова: ИНФОРМАЦИОННЫЕ ТЕХНОЛОГИИ, ИНФОРМАЦИОННЫЕ СИСТЕМЫ, ЦИФРОВАЯ СРЕДА, АВТОМАТИЗАЦИЯ, ОБРАБОТКА ДАННЫХ, ERP, CRM, BI, БАЗЫ ДАННЫХ, СУБД, ПОДДЕРЖКА ПРИНЯТИЯ РЕШЕНИЙ, КЛАССИФИКАЦИЯ ИС, БИЗНЕС-ПРОЦЕССЫ, ТЕХНИЧЕСКИЕ СРЕДСТВА, ПРОГРАММНОЕ ОБЕСПЕЧЕНИЕ, ТЕЛЕКОММУНИКАЦИИ, СТРУКТУРА ИС, УПРАВЛЕНИЕ, OLTP, DSS.

В работе исследуются сущность и взаимосвязь ключевых элементов цифровой среды — информационных технологий (ИТ) и информационных систем (ИС). Рассмотрены состав и компоненты современных ИТ, архитектурные принципы и функциональное назначение ИС, а также их классификация по различным критериям. Проанализирована трансформационная роль ИС в процессах управления данными и поддержки принятия решений в организациях.

Основные выводы:

\begin{enumerate}
\item Информационные технологии, как инструментарий, и информационные системы, как организационно-управленческие решения, образуют неразрывный симбиоз, лежащий в основе цифровой трансформации всех сфер общества.

\item Структура современной ИС, базирующаяся на чётком разделении уровней (данных, логики, представления), прямо определяет её функциональность, масштабируемость и эффективность в решении бизнес-задач.

\item Классификация ИС по сфере применения и уровню автоматизации (операционные, управленческие, стратегические) позволяет целенаправленно выбирать и внедрять решения, адекватные потребностям конкретной организации.

\item Эволюция ИС от простых средств автоматизации рутинных операций к сложным системам бизнес-аналитики (BI) и поддержки принятия решений (DSS) коренным образом меняет подходы к управлению, делая его более обоснованным и data-driven.
\end{enumerate}

Рекомендации:

\begin{itemize}[label=-]
\item При проектировании цифровой инфраструктуры организации рассматривать ИТ и ИС как единый комплекс, где технологии служат реализации системных целей.
\item Внедрять ИС, соответствующие уровню зрелости бизнес-процессов организации — от автоматизации транзакций (OLTP) до аналитической обработки (OLAP) и интеллектуального анализа данных.
\item Формировать корпоративные хранилища данных (Data Warehouse) как фундамент для эффективного управления информационными ресурсами и аналитики.
\item Использовать возможности современных ИС (ERP, CRM, BI-платформы) для создания единого информационного пространства и снижения рисков при принятии управленческих решений.
\end{itemize}

\newpage

\thispagestyle{empty}
\begin{center}
\bfseries СОДЕРЖАНИЕ
\end{center}
\vspace{-6pt}
\tableofcontents
\clearpage

% фантомные боли для секций, чтобы не было нумераций кроме тела реферата
\section*{}
\vspace*{-2.5em}
\begin{center}\textbf{ВВЕДЕНИЕ}\end{center}
\phantomsection
\addcontentsline{toc}{section}{Введение}

Информационные технологии и системы перестали быть лишь вспомогательным инструментом, превратившись в стратегический актив и основу функционирования современного общества — от государственного управления и бизнеса до повседневной жизни граждан. Их эффективное использование определяет конкурентоспособность, скорость адаптации к изменениям и качество принимаемых решений.

Современное состояние проблемы. Динамичное развитие ИТ (облачные вычисления, большие данные, IoT) постоянно расширяет возможности и усложняет архитектуру ИС. При этом существует разрыв между технологическим потенциалом и его реализацией в эффективных, безопасных и ориентированных на пользователя системах. Ключевой задачей становится не просто внедрение новейших технологий, а их грамотная интеграция в целостные ИС, решающие конкретные бизнес-задачи.

Актуальность темы обусловлена всеобъемлющим характером цифровизации, острой необходимостью в системном понимании взаимосвязи ИТ и ИС для преодоления технологической фрагментации, а также возрастающей ролью данных как ключевого ресурса развития.

Новизна исследования заключается в комплексном рассмотрении ИТ и ИС как двух сторон единого цифрового процесса, с акцентом на практическое значение их классификации и эволюционной роли в управлении и принятии решений.

Цель работы: исследовать понятийный аппарат, структуру, взаимосвязь и классификацию информационных технологий и систем, а также проанализировать их преобразующую роль в управлении данными и поддержке принятия решений.

Задачи исследования: раскрыть понятие и состав информационных технологий как элемента цифровой среды; проанализировать структуру и функции информационных систем в современном обществе; исследовать взаимосвязь и взаимозависимость между информационными технологиями и информационными системами;
провести классификацию информационных систем по сфере применения и уровню автоматизации; оценить роль информационных систем в процессах управления данными и поддержки принятия решений.

Методы исследования: системный и структурно-функциональныйанализ для изучения ИТ и ИС как целостных комплексов; классификация и сравнительный анализ для категоризации информационных систем по различным критериям; анализ литературных и нормативных источников (ГОСТ, отраслевые стандарты) для определения базовых понятий; анализ практических кейсов внедрения ИС для выявления их роли в управленческих процессах.

\newpage
% вручную увеличить номер секции и записать в оглавление с номером,
% чтобы в теле не печатался автоматический номер, но в tableofcontents он остался
\refstepcounter{section}
\addcontentsline{toc}{section}{\protect\numberline{\thesection} Понятие и состав информационных технологий как элемента цифровой среды}
\vspace*{-2.5em}
\begin{center}\textbf{\thesection\ ПОНЯТИЕ И СОСТАВ ИНФОРМАЦИОННЫХ ТЕХНОЛОГИЙ КАК ЭЛЕМЕНТА ЦИФРОВОЙ СРЕДЫ}\end{center}
\vspace*{-1.5em}
\ManualSubsection{Понятие информационных технологий}

Информационные технологии (ИТ), часто расширяемые до понятия «информационно-коммуникационные технологии» (ИКТ), представляют собой комплекс взаимосвязанных научных, технологических и инженерных дисциплин, изучающих методы эффективной организации труда людей, занятых обработкой и хранением информации, а также вычислительную технику и методы организации и взаимодействия с людьми и производственным оборудованием \cite{rosstat_it_def}. В современном контексте ИТ выступают фундаментальным элементом цифровой среды — технологической основы, на которой строится функционирование общества, экономики и государственного управления. Цифровая среда формируется как результат интеграции и синергетического взаимодействия разнородных информационных технологий, создавая новое качество информационного пространства \cite{habr_digital_env}.

\ManualSubsection{Аппаратное обеспечение}

Аппаратное обеспечение составляет физическую основу ИТ, представляя собой совокупность электронных и механических устройств. К этому компоненту относятся не только конечные пользовательские устройства, но и инфраструктурное оборудование: серверные системы, обеспечивающие хранение и обработку данных; сетевое оборудование, формирующее каналы передачи информации; а также периферийные устройства для ввода и вывода данных. Данный слой является материальным носителем для выполнения вычислительных операций.

\ManualSubsection{Программное обеспечение}

Программное обеспечение является логической надстройкой, определяющей алгоритмы функционирования аппаратных средств. Оно включает системное ПО, в первую очередь операционные системы, которые управляют ресурсами и предоставляют базовые сервисы. Прикладное ПО решает конкретные задачи пользователя или организации — от офисных пакетов до сложных корпоративных систем. Качество и архитектура программного слоя напрямую определяют функциональность и безопасность всей информационной системы \cite{techopedia_it_components}.

\ManualSubsection{Коммуникационные технологии}

Коммуникационный компонент обеспечивает интеграцию разрозненных аппаратных и программных ресурсов в единую информационную среду. Его основу составляют компьютерные сети различных масштабов — от локальных до глобальных. Данный сегмент включает как физические среды передачи данных, так и протоколы связи, которые регламентируют порядок взаимодействия. Именно сетевые технологии трансформируют совокупность отдельных компьютеров в распределённую систему.

\ManualSubsection{Информационные ресурсы}

Данные представляют собой формализованную информацию, пригодную для обработки автоматическими средствами. Этот компонент является семантическим ядром ИТ, поскольку вся технологическая инфраструктура создаётся именно для работы с данными. Информационные ресурсы могут быть структурированы, слабоструктурированы или неструктурированы. Управление данными, включающее процессы их создания, хранения, актуализации и анализа, является одной из центральных задач.

\ManualSubsection{Кадровый и организационный компонент}

Технологии не существуют в вакууме; их разработка, внедрение и эксплуатация осуществляются людьми и регулируются определёнными правилами. Кадровый компонент охватывает специалистов различного профиля. Организационно-управленческий аспект включает стандарты, регламенты, методики и правовые нормы, которые определяют, как технологии должны применяться для достижения целей \cite{habr_digital_env}.

\ManualSubsection{Синтез компонентов}

Целостность цифровой среды достигается сложным взаимодействием всех описанных компонентов. Аппаратное обеспечение предоставляет платформу, программное — алгоритмы, сети обеспечивают коммуникацию, данные наполняют систему смыслом, а человеческий фактор задаёт цели. Например, процесс онлайн-платежа представляет собой скоординированную последовательность действий всех компонентов. Только в таком синтезе информационные технологии реализуют свою основную функцию — превращение данных в знания и решения \cite{techopedia_it_components}.

\newpage
\refstepcounter{section}
\addcontentsline{toc}{section}{\protect\numberline{\thesection}Структура и функции информационных систем в современном обществе}
\begin{center}\textbf{\thesection\ СТРУКТУРА И ФУНКЦИИ ИНФОРМАЦИОННЫХ СИСТЕМ В СОВРЕМЕННОМ ОБЩЕСТВЕ}\end{center}
\vspace*{-1.5em}
\ManualSubsection{Понятие и сущность информационной системы}

Информационная система (ИС) представляет собой организационно-технический комплекс, предназначенный для сбора, хранения, обработки, анализа и распространения информации для решения управленческих и операционных задач. Это система, компоненты которой включают данные, технические и программные средства, а также персонал, объединенные для достижения общей цели \cite{stair_reynolds_2018}. В современном контексте ИС трансформируются из инструментов автоматизации рутинных операций в стратегические активы, формирующие цифровую экосистему общества и определяющие его эффективность и конкурентоспособность \cite{oecd_digital_economy_2019}.

\ManualSubsection{Структурные компоненты информационных систем}

Фундаментальную основу ИС составляют данные, выступающие объектом обработки и хранения. Данные, структурированные в соответствии с моделью предметной области, формируют информационную базу системы. Технические средства (аппаратное обеспечение) образуют физическую платформу ИС, включая серверное оборудование, устройства хранения, сетевую инфраструктуру и клиентские рабочие станции. Программное обеспечение определяет логику функционирования системы, реализуя алгоритмы обработки данных и предоставляя интерфейсы для взаимодействия. Критически важным компонентом является персонал, который подразделяется на пользователей, эксплуатирующих систему для решения прикладных задач, и технических специалистов, обеспечивающих её разработку, внедрение и сопровождение. Организационно-правовая компонента, включающая регламенты, стандарты и политики, формализует порядок использования ИС и управления её жизненным циклом.

\ManualSubsection{Функциональный спектр информационных систем}

Первичной функцией ИС является сбор и регистрация данных из внутренних и внешних источников. Функция обработки и преобразования данных включает их структурирование, агрегацию, вычисление и анализ, трансформируя исходные данные в значимую информацию. Организация надёжного и безопасного хранения информации, обеспечивающая её целостность, доступность и конфиденциальность, представляет собой ключевую задачу, решаемую современными системами управления базами данных и хранилищами \cite{nist_cybersecurity_framework_2018}. Функция поиска и предоставления информации реализуется через различные интерфейсы и средства генерации отчётов, обеспечивая пользователей актуальными сведениями. Поддержка принятия решений осуществляется за счёт аналитических и моделирующих подсистем, которые предоставляют инструменты для оценки альтернатив и прогнозирования. Функция интеграции обеспечивает взаимодействие с другими информационными системами, создавая единое информационное пространство организации или отрасли.

\ManualSubsection{Роль информационных систем в обществе}

В экономической сфере ИС составляют основу цифровой трансформации бизнеса, оптимизируя цепочки создания стоимости, обеспечивая управление ресурсами предприятия (ERP) и взаимоотношениями с клиентами (CRM). В системе государственного управления ИС реализуют концепцию «электронного правительства», повышая прозрачность, эффективность предоставления услуг и качество управленческих решений. В социальной сфере, включая здравоохранение и образование, информационные системы способствуют персонализации услуг, удалённому доступу и повышению общего качества жизни. В научно-исследовательской деятельности ИС обеспечивают обработку больших данных, вычислительные эксперименты и глобальную коллаборацию учёных. Таким образом, информационные системы превратились в критическую инфраструктуру современного общества, определяя траектории его развития и формируя новые вызовы в области кибербезопасности, цифрового суверенитета и этики использования данных \cite{oecd_digital_economy_2019}.

\newpage
\refstepcounter{section}
\addcontentsline{toc}{section}{\protect\numberline{\thesection}Взаимосвязь между информационными технологиями и информационными системами}
\vspace*{-1.5em}
\begin{center}\textbf{\thesection\ ВЗАИМОСВЯЗЬ МЕЖДУ ИНФОРМАЦИОННЫМИ ТЕХНОЛОГИЯМИ И ИНФОРМАЦИОННЫМИ СИСТЕМАМИ}\end{center}
\vspace*{-1.5em}

\ManualSubsection{Концептуальные основы взаимосвязи}

Информационные технологии (ИТ) и информационные системы (ИС) представляют собой взаимозависимые, но концептуально различные категории. Информационные технологии определяются как совокупность методов, производственных процессов и программно-технических средств, объединенных в технологическую цепочку, обеспечивающую сбор, хранение, обработку, вывод и распространение информации \cite{it-foundation}. В свою очередь, информационная система — это комплекс взаимосвязанных компонентов, который собирает, обрабатывает, хранит и распределяет информацию для поддержки принятия решений, координации, контроля, анализа и визуализации в организации \cite{is-definition}. Таким образом, ИТ служат инструментальной основой, в то время как ИС представляют собой целостное организационно-управленческое решение, использующее эти инструменты для достижения конкретных целей.

\ManualSubsection{ИТ как технологический фундамент ИС}

Любая информационная система базируется на определённом наборе информационных технологий. Аппаратное обеспечение, включая серверы, устройства хранения и сетевую инфраструктуру, формирует физическую платформу системы. Программное обеспечение, от операционных систем и систем управления базами данных до специализированных прикладных программ, реализует алгоритмы обработки данных и предоставляет функциональные возможности. Телекоммуникационные технологии обеспечивают интеграцию компонентов и взаимодействие с пользователями. Без этого технологического стека существование современной сложной ИС невозможно, что подтверждает тезис о фундаментальной зависимости систем от развитости технологий \cite{systems-approach}. Например, система онлайн-банкинга является практической реализацией таких технологий, как веб-сервисы, криптография и реляционные СУБД.

\ManualSubsection{ИС как целевой контекст применения ИТ}

Обратное влияние заключается в том, что именно потребности и задачи информационных систем выступают основным драйвером развития информационных технологий. Необходимость решения новых классов задач (например, обработки больших данных, реализации сложных бизнес-процессов или обеспечения киберустойчивости) стимулирует исследования и разработки в технологической сфере. ИС задают требования к производительности, надёжности, безопасности и масштабируемости, которым должны соответствовать ИТ. Этот симбиоз формирует цикл развития: усовершенствованные технологии открывают возможность для создания более эффективных и функциональных систем, а растущие амбиции систем проектируют дорожную карту для эволюции технологий \cite{systems-approach}.

\ManualSubsection{Практические модели взаимодействия}

На практике взаимосвязь проявляется в архитектурных моделях, где выбор конкретных технологий напрямую определяет характеристики будущей системы. Концепция N-уровневой архитектуры явно разделяет технологические слои (уровень данных, уровень бизнес-логики, уровень представления), но интегрирует их в единую систему. Развитие облачных технологий (IaaS, PaaS, SaaS) трансформирует подход к построению ИС, смещая акцент с владения инфраструктурой на управление сервисами. Другим примером является внедрение парадигмы микросервисов, где технология контейнеризации (Docker, Kubernetes) позволяет декомпозировать монолитную систему на набор слабосвязанных, независимо развертываемых сервисов, повышая гибкость и отказоустойчивость всей ИС в целом.

\ManualSubsection{Интеграция и синергетический эффект}

Конечная ценность создаётся не технологиями или системами по отдельности, а их эффективной интеграцией. Успешная информационная система представляет собой синергетическое целое, где технологические компоненты подобраны и настроены для оптимального выполнения бизнес-функций. Эта интеграция требует системного подхода, учитывающего не только техническую совместимость, но и организационные аспекты, такие как соответствие бизнес-процессам и квалификация персонала. Таким образом, взаимосвязь ИТ и ИС имеет диалектический характер: технологии материализуют системные замыслы, а системы придают технологиям практическую значимость и направление для развития, совместно формируя цифровой ландшафт современного общества \cite{is-definition}.

\newpage
\refstepcounter{section}
\addcontentsline{toc}{section}{\protect\numberline{\thesection}Классификация информационных систем по сфере применения и уровню автоматизации}
\begin{center}\textbf{\thesection\ КЛАССИФИКАЦИЯ ИНФОРМАЦИОННЫХ СИСТЕМ ПО СФЕРЕ ПРИМЕНЕНИЯ И УРОВНЮ АВТОМАТИЗАЦИИ}\end{center}
\vspace*{-1.5em}
\ManualSubsection{Методологические основы классификации}

Классификация информационных систем является необходимым инструментом для их анализа, сравнения и выбора, позволяя структурировать обширное множество разнородных решений по значимым критериям. В основе систематизации лежит понимание ИС как артефакта, созданного для поддержки определённых процессов в рамках заданного организационного или социального контекста \cite{iso-iec-15288}. Наиболее релевантными для практического применения признаны два ключевых критерия: целевая сфера деятельности, которую система обслуживает, и степень автоматизации функций, которые она выполняет. Данные критерии отражают как внешние требования к системе, так и её внутреннюю зрелость.

\ManualSubsection{Классификация по сфере применения}

Данный подход группирует системы в соответствии с предметной областью и специфическими бизнес-задачами. В коммерческом секторе широкое распространение получили системы управления взаимоотношениями с клиентами (CRM) и автоматизации продаж, направленные на повышение лояльности и увеличение доходов. Управление ресурсами предприятия (ERP) представляет собой интегрированные платформы, объединяющие финансовый учёт, логистику, управление персоналом и производственные операции в единый контур данных. Для государственного сектора характерны системы электронного документооборота (СЭД) и порталы государственных услуг, основными задачами которых являются повышение прозрачности, доступности сервисов и эффективности внутренних административных процессов. В социальной сфере, такой как здравоохранение и образование, используются специализированные ИС: электронные медицинские карты (EMR/EHR) и системы управления обучением (LMS), которые фокусируются на персонализированном обслуживании и управлении знаниями \cite{oecd-digital-economy}. Эта классификация помогает организациям идентифицировать класс систем, наиболее адекватный их отраслевой специфике.

\ManualSubsection{Классификация по уровню автоматизации}

Уровень автоматизации определяет степень, в которой система замещает ручной человеческий труд и реализует сложную логику принятия решений. На нижнем уровне находятся транзакционные системы обработки операций (OLTP), основная функция которых заключается в регистрации, хранении и первичной обработке структурированных данных, таких как продажи или платежи. Следующий уровень представлен информационными системами управления (MIS), которые агрегируют данные из OLTP-систем для формирования стандартизированных отчётов о текущей деятельности, обеспечивая операционный контроль. Более высокую ступень занимают системы поддержки принятия решений (DSS), которые предоставляют аналитические инструменты, моделирование сценариев и интерактивный анализ данных для помощи менеджерам в решении слабоструктурированных задач. На вершине эволюционной пирамиды находятся экспертные системы (ES) и системы на основе искусственного интеллекта, способные к автоматическому анализу сложных ситуаций, выработке рекомендаций и, в некоторых случаях, автономному выполнению действий в соответствии с заложенными знаниями и алгоритмами \cite{iso-iec-15288}.

\ManualSubsection{Взаимодействие критериев и гибридные системы}

На практике сфера применения и уровень автоматизации тесно переплетаются, формируя многомерное пространство классификации. Например, в производственной сфере система SCADA (диспетчерское управление и сбор данных) является одновременно отраслевым решением и системой с высоким уровнем автоматизации технологических процессов. Современная ERP-система интегрирует в себе функции транзакционной обработки (учёт), информационного управления (отчётность) и поддержки принятия решений (планирование ресурсов). Развитие технологий, таких как большие данные и машинное обучение, приводит к появлению гибридных систем, которые сложно однозначно отнести к одному классу. Система рекомендаций в электронной коммерции, будучи отраслевым решением для ритейла, использует алгоритмы искусственного интеллекта, характерные для наиболее высокого уровня автоматизации анализа \cite{gartner-it-glossary}. Таким образом, актуальная классификация должна рассматриваться как динамическая система координат, а не как набор жёстких категорий.

\newpage
\refstepcounter{section}
\addcontentsline{toc}{section}{\protect\numberline{\thesection}Роль информационных систем в управлении данными и поддержке принятия решений}
\begin{center}\textbf{\thesection\ РОЛЬ ИНФОРМАЦИОННЫХ СИСТЕМ В УПРАВЛЕНИИ ДАННЫМИ И ПОДДЕРЖКЕ ПРИНЯТИЯ РЕШЕНИЙ}\end{center}
\vspace*{-1.5em}
\ManualSubsection{Информационные системы как основа управления данными}

Информационные системы являются центральным технологическим звеном в современном управлении корпоративными и общественными данными. Они реализуют полный жизненный цикл данных, начиная с автоматизированного сбора из разнородных источников (датчиков, транзакционных систем, внешних сервисов), обеспечения их качества посредством валидации и очистки, и заканчивая организацией структурированного хранения в системах управления базами данных (СУБД) или хранилищах данных. Ключевая роль ИС заключается в трансформации сырых данных в качественный информационный актив, доступный для дальнейшего использования. Эффективное управление данными через ИС подразумевает не только их сохранность и целостность, но и обеспечение безопасности, регулируемого доступа, а также поддержание актуальности, что формирует надёжный фундамент для любой аналитической деятельности \cite{is_management}.

\ManualSubsection{От данных к решениям: аналитический контур ИС}

Современные информационные системы вышли за рамки пассивного хранения информации, сформировав активный аналитический контур. Этот контур представлен классами систем, специально предназначенных для поддержки принятия решений (Decision Support Systems, DSS). В отличие от операционных систем, DSS фокусируются на анализе агрегированных исторических и текущих данных, использовании математических моделей и методов имитационного моделирования для оценки различных сценариев развития ситуации. Они предоставляют лицам, принимающим решения, не просто отчёты, а интерактивные инструменты для «проигрывания» вариантов, выявления скрытых зависимостей и оценки потенциальных рисков и последствий \cite{dss_overview}. Таким образом, ИС выступает в роли интеллектуального посредника, сокращающего когнитивную нагрузку на управленца и переводящего процесс принятия решений с интуитивного на доказательный уровень.

\ManualSubsection{Иерархия решений и соответствующие классы ИС}

Поддержка, оказываемая информационными системами, дифференцируется в соответствии с иерархией управленческих решений. Для оперативного уровня, связанного с ежедневной рутинной деятельностью, транзакционные системы (OLTP) обеспечивают основу, предоставляя актуальные данные о текущем состоянии процессов (например, об остатках на складе). На тактическом уровне, где решения принимаются на среднесрочную перспективу (например, планирование бюджета или маркетинговой кампании), используются системы управления эффективностью бизнеса и бизнес-аналитики (BI), которые агрегируют и анализируют данные за определённые периоды. Стратегические решения, определяющие долгосрочное развитие организации, требуют наиболее сложных систем, включающих прогнозное моделирование, анализ больших данных и элементы искусственного интеллекта для выявления глобальных тенденций и оценки рыночных возможностей \cite{dss_overview}. Эта специализация демонстрирует, как архитектура ИС-ландшафта организации должна соответствовать структуре её управленческих задач.

\ManualSubsection{Принципы data-driven управления и роль ИС}

Доминирующей парадигмой современного менеджмента становится управление на основе данных (data-driven management), и информационные системы являются его технологическим воплощением. Данный подход утверждает, что ключевые управленческие решения должны приниматься не на основе интуиции или прошлого опыта, а на основании объективного анализа количественных и качественных данных. ИС делает этот принцип реализуемым, обеспечивая сбор релевантных метрик, их визуализацию в понятных дашбордах и предоставляя инструменты для глубокой аналитики. Это смещает культуру управления в сторону большей прозрачности, измеримости и accountability \cite{data_driven}. Успешные организации используют ИС для создания единого источника истины о бизнесе, что позволяет согласованно действовать разным подразделениям, основываясь на одних и тех же фактах.

\ManualSubsection{Эволюция и перспективы: от отчётности к интеллектуальной аналитике}

Роль информационных систем в поддержке решений продолжает эволюционировать. Если первоначально они автоматизировали формирование стандартных отчётов, то сегодня фокус смещается на предиктивную и прескриптивную аналитику. Современные системы способны не только описывать, что произошло (дескриптивная аналитика), но и прогнозировать, что может произойти (предиктивная аналитика), а также предлагать конкретные действия для достижения желаемых результатов (прескриптивная аналитика). Интеграция технологий машинного обучения и обработки естественного языка позволяет создавать интеллектуальных помощников и системы рекомендаций нового поколения. Таким образом, информационные системы трансформируются из инструмента реагирования в стратегическую платформу для проактивного управления, становясь ключевым фактором устойчивого развития и конкурентного преимущества в цифровую эпоху \cite{data_driven}.

\newpage
\section*{}
\vspace*{-2.5em}
\begin{center}\textbf{ЗАКЛЮЧЕНИЕ}\end{center}
\phantomsection
\addcontentsline{toc}{section}{Заключение}

Информационные технологии и системы образуют фундаментальный симбиоз, лежащий в основе цифровой трансформации всех сфер современного общества. В ходе проведённого исследования были раскрыты состав и структура ИТ как элемента цифровой среды, определены архитектурные принципы и функциональное назначение информационных систем, а также проанализирована их эволюционная роль в процессах управления данными и поддержки принятия решений.

Установлено, что информационные технологии, включающие аппаратное и программное обеспечение, средства связи и организационные компоненты, служат инструментальной базой. В свою очередь, информационные системы представляют собой целостные организационно-управленческие решения, использующие этот инструментарий для достижения конкретных бизнес-целей. Классификация ИС по сфере применения и уровню автоматизации демонстрирует их адаптивность и специализацию, позволяя организациям выбирать решения, адекватные их операционным и стратегическим потребностям.

Ключевым выводом является трансформация роли информационных систем от средств автоматизации рутинных операций до стратегических платформ, обеспечивающих управление на основе данных (data-driven management). Современные ИС, особенно системы поддержки принятия решений (DSS) и бизнес-аналитики (BI), превращают данные в ценную информацию, прогнозы и рекомендации, существенно повышая качество и обоснованность управленческих решений. Таким образом, эффективная интеграция информационных технологий и систем становится критическим фактором конкурентоспособности, устойчивого развития и формирования цифровой экосистемы будущего.

\newpage
\begin{center}\textbf{СПИСОК СОКРАЩЕНИЙ}\end{center}
\phantomsection
\addcontentsline{toc}{section}{Список сокращений}

\vspace{1em}

{
\small
\noindent
\begin{tabular}{|l|p{0.75\textwidth}|}
\hline
\textbf{Сокращение} & \multicolumn{1}{c|}{\textbf{Расшифровка}} \\
\noalign{\hrule height 1.2pt}
\hline
ИТ & Информационные технологии. Совокупность методов, процессов и программно-технических средств для работы с информацией. \\
\hline
ИКТ & Информационно-коммуникационные технологии. Расширенное понятие ИТ, включающее телекоммуникационные аспекты.\\
\hline
ИС & Информационная система. Организационно-технический комплекс для сбора, хранения, обработки и распространения информации.\\
\hline
Цифровая среда & Технологическая и социальная экосистема, формируемая взаимодействием ИТ, ИС, данных, людей и процессов.\\
\hline
СУБД & Система управления базами данных. Программное обеспечение для создания, управления и использования баз данных. \\
\hline
ERP & Планирование ресурсов предприятия. Интегрированная система для управления ключевыми бизнес-процессами компании. \\
\hline
CRM & Управление взаимоотношениями с клиентами. Система для автоматизации процессов взаимодействия с клиентами. \\
\hline
BI & Бизнес-аналитика. Технологии и практики для анализа данных и предоставления информации, полезной для принятия решений.\\
\hline
DSS & Система поддержки принятия решений. Класс ИС, предназначенный для помощи в решении слабоструктурированных управленческих задач. \\
\hline
LMS & Система управления обучением. Программная платформа для администрирования, документирования и предоставления образовательных курсов.\\
\hline
OLTP & Обработка транзакций в реальном времени. Класс систем, ориентированных на выполнение и регистрацию повседневных операций.\\
\hline
NoSQL & Not Only SQL. Подход к проектированию баз данных, отличный от классических реляционных моделей.\\
\hline
API & Интерфейс программирования приложений. Набор определений и протоколов для взаимодействия между программными компонентами. \\
\hline
\end{tabular}
}

\clearpage
\nocite{*}
\printbibliography[heading=bibliography]

\end{document}
